% GENERATED FILE. Edit canonical lessons, then run npm run generate:books.
% canonical-source-hash: fnv1a64:1372ad0d41f07cd7
% canonical-lessons: RU-C310-yubiley, RU-C310-godovshchina, RU-C310-novorozhdyonnyy, RU-C310-krestiny, RU-C310-krestit, RU-R310-review-303-310

\chapter{Anniversaries and christenings}
\label{ch:ru-anniversaries-and-christenings}
\hlchaptermodality{310}

\begin{chapteropening}
\textbf{By the end of this chapter:} \emph{I can talk about a big birthday, an anniversary, a newborn, a christening party and christening a child.}

Complete the last lesson of chapter 310: I can talk about a big birthday, an anniversary, a newborn, a christening party and christening a child.
\end{chapteropening}

\section[\texorpdfstring{\ru{юбилей} (yubiléy)}{юбилей (yubiléy)}]{\ru{юбилей} (yubiléy) — a jubilee, a big round birthday or anniversary}

\label{lesson:RU-C310-yubiley}



\begin{quote}
Before the new one: say the Russian for a widow, then the Russian for an heir.
\end{quote}

\subsection*{You'll want to know: \ru{юбилей}}
\textbf{\ru{юбилей}} (\emph{yubiléy}) — \textquotedblleft{}a jubilee, a big round birthday or anniversary\textquotedblright{}.

\begin{quote}
\ru{У} \ru{дедушки} \ru{скоро} \ru{юбилей}, \ru{семьдесят} \ru{лет}.
\end{quote}

\begin{quote}
\emph{U dédushki skóro yubiléy, sémdesyat let.}
\end{quote}

\begin{quote}
Grandpa soon has a big birthday: seventy.
\end{quote}

A round number: fifty, sixty, seventy years.

\begin{cousinweb}
\emph{\ru{юбилей}} was borrowed from German \emph{Jubiläum} (EtymDB), which comes from Late Latin \emph{iūbilaeus} \textquotedblleft{}the jubilee year\textquotedblright{}, which is usually traced to a Hebrew word for the ram's horn that announced it.
\end{cousinweb}

\begin{grammarlens}[title={What you've built}]
One more word for telling, explaining and describing at B1.
\end{grammarlens}

\subsection*{Your turn}
\begin{itemize}
\raggedright
  \item \emph{Say it:} \emph{yubiléy}
  \item \emph{Say it:} \emph{yubiléy}, once more
  \item \emph{Say it:} a sentence of your own with \emph{yubiléy}
  \item \emph{Recall:} say the Russian for a widow, then the Russian for an heir, then say \emph{yubiléy} again
\end{itemize}

\begin{tcolorbox}[breakable,colback=teal!4,colframe=teal!35!black,title={Before you move on}]
What does \ru{юбилей} mean? (\textquotedblleft{}A jubilee, a big round birthday or anniversary\textquotedblright{}.) Say it once more.
\end{tcolorbox}
\section[\texorpdfstring{\ru{годовщина} (godovshchína)}{годовщина (godovshchína)}]{\ru{годовщина} (godovshchína) — an anniversary}

\label{lesson:RU-C310-godovshchina}



\begin{quote}
Before the new one: say the Russian for an heir, then the Russian for a jubilee.
\end{quote}

\subsection*{You'll want to know: \ru{годовщина}}
\textbf{\ru{годовщина}} (\emph{godovshchína}) — \textquotedblleft{}an anniversary\textquotedblright{}.

\begin{quote}
\ru{Сегодня} \ru{годовщина} \ru{их} \ru{свадьбы}.
\end{quote}

\begin{quote}
\emph{Sevódnya godovshchína ikh svádby.}
\end{quote}

\begin{quote}
Today is their wedding anniversary.
\end{quote}

Feminine.

\begin{cousinweb}
\emph{\ru{годовщина}} is built on \emph{\ru{год}} \textquotedblleft{}year\textquotedblright{}, inherited through Old East Slavic \emph{godŭ} from Proto-Slavic \emph{godŭ} \textquotedblleft{}the right time\textquotedblright{} (EtymDB).
\end{cousinweb}

\begin{grammarlens}[title={What you've built}]
One more word for telling, explaining and describing at B1.
\end{grammarlens}

\subsection*{Your turn}
\begin{itemize}
\raggedright
  \item \emph{Say it:} \emph{godovshchína}
  \item \emph{Say it:} \emph{godovshchína}, once more
  \item \emph{Say it:} a sentence of your own with \emph{godovshchína}
  \item \emph{Recall:} say the Russian for an heir, then the Russian for a jubilee, then say \emph{godovshchína} again
\end{itemize}

\begin{tcolorbox}[breakable,colback=teal!4,colframe=teal!35!black,title={Before you move on}]
What does \ru{годовщина} mean? (\textquotedblleft{}An anniversary\textquotedblright{}.) Say it once more.
\end{tcolorbox}
\section[\texorpdfstring{\ru{новорождённый} (novorozhdyónnyy)}{новорождённый (novorozhdyónnyy)}]{\ru{новорождённый} (novorozhdyónnyy) — a newborn}

\label{lesson:RU-C310-novorozhdyonnyy}



\begin{quote}
Before the new one: say the Russian for a jubilee, then the Russian for an anniversary.
\end{quote}

\subsection*{You'll want to know: \ru{новорождённый}}
\textbf{\ru{новорождённый}} (\emph{novorozhdyónnyy}) — \textquotedblleft{}a newborn\textquotedblright{}.

\begin{quote}
\ru{Новорождённый} \ru{весил} \ru{три} \ru{килограмма}.
\end{quote}

\begin{quote}
\emph{Novorozhdyónnyy vésil tri kilográmma.}
\end{quote}

\begin{quote}
The newborn weighed three kilograms.
\end{quote}

An adjective that serves as a noun. A girl is \emph{\ru{новорождённая}} (\emph{novorozhdyónnaya}).

\begin{cousinweb}
\emph{\ru{новорождённый}} joins \emph{\ru{новый}} \textquotedblleft{}new\textquotedblright{}, inherited from Proto-Slavic \emph{novŭ}, PIE \emph{néwos}, and \emph{\ru{рождённый}} \textquotedblleft{}born\textquotedblright{}, from \emph{\ru{родить}}, Proto-Slavic \emph{roditi} (EtymDB).
\end{cousinweb}

\begin{grammarlens}[title={What you've built}]
One more word for telling, explaining and describing at B1.
\end{grammarlens}

\subsection*{Your turn}
\begin{itemize}
\raggedright
  \item \emph{Say it:} \emph{novorozhdyónnyy}
  \item \emph{Say it:} \emph{novorozhdyónnyy}, once more
  \item \emph{Say it:} a sentence of your own with \emph{novorozhdyónnyy}
  \item \emph{Recall:} say the Russian for a jubilee, then the Russian for an anniversary, then say \emph{novorozhdyónnyy} again
\end{itemize}

\begin{tcolorbox}[breakable,colback=teal!4,colframe=teal!35!black,title={Before you move on}]
What does \ru{новорождённый} mean? (\textquotedblleft{}A newborn\textquotedblright{}.) Say it once more.
\end{tcolorbox}
\section[\texorpdfstring{\ru{крестины} (krestíny)}{крестины (krestíny)}]{\ru{крестины} (krestíny) — a christening (party)}

\label{lesson:RU-C310-krestiny}



\begin{quote}
Before the new one: say the Russian for an anniversary, then the Russian for a newborn.
\end{quote}

\subsection*{You'll want to know: \ru{крестины}}
\textbf{\ru{крестины}} (\emph{krestíny}) — \textquotedblleft{}a christening (party)\textquotedblright{}.

\begin{quote}
\ru{На} \ru{крестины} \ru{пригласили} \ru{всех} \ru{родственников}.
\end{quote}

\begin{quote}
\emph{Na krestíny priglasíli vsekh ródstvennikov.}
\end{quote}

\begin{quote}
All the relatives were invited to the christening.
\end{quote}

Plural only; genitive \emph{\ru{крестин}}.

\begin{cousinweb}
\emph{\ru{крестины}} is built on \emph{\ru{крест}} \textquotedblleft{}a cross\textquotedblright{}, from Church Slavonic \emph{krĭstŭ}, which EtymDB traces ultimately to Greek \emph{Khristós} \textquotedblleft{}the anointed\textquotedblright{}.
\end{cousinweb}

\begin{grammarlens}[title={What you've built}]
One more word for telling, explaining and describing at B1.
\end{grammarlens}

\subsection*{Your turn}
\begin{itemize}
\raggedright
  \item \emph{Say it:} \emph{krestíny}
  \item \emph{Say it:} \emph{krestíny}, once more
  \item \emph{Say it:} a sentence of your own with \emph{krestíny}
  \item \emph{Recall:} say the Russian for an anniversary, then the Russian for a newborn, then say \emph{krestíny} again
\end{itemize}

\begin{tcolorbox}[breakable,colback=teal!4,colframe=teal!35!black,title={Before you move on}]
What does \ru{крестины} mean? (\textquotedblleft{}A christening (party)\textquotedblright{}.) Say it once more.
\end{tcolorbox}
\section[\texorpdfstring{\ru{крестить} (krestít)}{крестить (krestít)}]{\ru{крестить} (krestít) — to christen, to baptize}

\label{lesson:RU-C310-krestit}



\begin{quote}
Before the new one: say the Russian for a newborn, then the Russian for a christening.
\end{quote}

\subsection*{You'll want to know: \ru{крестить}}
\textbf{\ru{крестить}} (\emph{krestít}) — \textquotedblleft{}to christen, to baptize\textquotedblright{}.

\begin{quote}
\ru{Ребёнка} \ru{крестили} \ru{в} \ru{мае}.
\end{quote}

\begin{quote}
\emph{Rebyónka krestíli v máye.}
\end{quote}

\begin{quote}
The baby was christened in May.
\end{quote}

Both aspects; also perfective \emph{\ru{окрестить}} (\emph{okrestít}). Present \emph{\ru{крещу}}, \emph{\ru{крестишь}}.

\begin{cousinweb}
\emph{\ru{крестить}} is inherited from Proto-Slavic \emph{krĭstiti}, built on \emph{krĭstŭ} \textquotedblleft{}baptism\textquotedblright{} (EtymDB), a word that goes back to Greek \emph{Khristós}.
\end{cousinweb}

\begin{grammarlens}[title={What you've built}]
One more word for telling, explaining and describing at B1.
\end{grammarlens}

\subsection*{Your turn}
\begin{itemize}
\raggedright
  \item \emph{Say it:} \emph{krestít}
  \item \emph{Say it:} \emph{krestít}, once more
  \item \emph{Say it:} a sentence of your own with \emph{krestít}
  \item \emph{Recall:} say the Russian for a newborn, then the Russian for a christening, then say \emph{krestít} again
\end{itemize}

\begin{tcolorbox}[breakable,colback=teal!4,colframe=teal!35!black,title={Before you move on}]
What does \ru{крестить} mean? (\textquotedblleft{}To christen, to baptize\textquotedblright{}.) Say it once more.
\end{tcolorbox}
\section[(dialogue)]{Words from chapters 303–310 — a review}

\label{lesson:RU-R310-review-303-310}



\begin{quote}
Say the Russian for birth and to be born.
\end{quote}

\subsection*{Your turn}
Say these aloud:

\begin{itemize}
\raggedright
  \item birth, to be born, a baby, pregnant, to give birth
  \item childhood, youth, one's young years, old age, to grow up
  \item to bring up, to adopt, to grow old, a graduation party, a generation
  \item an engagement, to fall in love, to court, a church wedding, to part
  \item a divorce, to get divorced, a move, to quit, an inheritance
  \item a school-leaver, to enter, to finish, a degree, a school-leaving certificate
  \item a funeral, to bury, death, a widow, an heir
  \item a jubilee, an anniversary, a newborn, a christening, to christen
\end{itemize}

\begin{tcolorbox}[breakable,colback=teal!4,colframe=teal!35!black,title={Before you move on}]
Birth? (\textbf{\ru{рождение}}.) To be born? (\textbf{\ru{рождаться}}.) A baby? (\textbf{\ru{младенец}}.) Pregnant? (\textbf{\ru{беременная}}.) To give birth? (\textbf{\ru{рожать}}.) Childhood? (\textbf{\ru{детство}}.) Youth? (\textbf{\ru{юность}}.) One's young years? (\textbf{\ru{молодость}}.) Old age? (\textbf{\ru{старость}}.) To grow up? (\textbf{\ru{взрослеть}}.) To bring up? (\textbf{\ru{воспитывать}}.) To adopt? (\textbf{\ru{усыновлять}}.) To grow old? (\textbf{\ru{стареть}}.) A graduation party? (\textbf{\ru{выпускной}}.) A generation? (\textbf{\ru{поколение}}.) An engagement? (\textbf{\ru{помолвка}}.) To fall in love? (\textbf{\ru{влюбляться}}.) To court? (\textbf{\ru{ухаживать}}.) A church wedding? (\textbf{\ru{венчание}}.) To part? (\textbf{\ru{расставаться}}.) A divorce? (\textbf{\ru{развод}}.) To get divorced? (\textbf{\ru{разводиться}}.) A move? (\textbf{\ru{переезд}}.) To quit? (\textbf{\ru{увольняться}}.) An inheritance? (\textbf{\ru{наследство}}.) A school-leaver? (\textbf{\ru{выпускник}}.) To enter? (\textbf{\ru{поступать}}.) To finish? (\textbf{\ru{оканчивать}}.) A degree? (\textbf{\ru{диплом}}.) A school-leaving certificate? (\textbf{\ru{аттестат}}.) A funeral? (\textbf{\ru{похороны}}.) To bury? (\textbf{\ru{хоронить}}.) Death? (\textbf{\ru{смерть}}.) A widow? (\textbf{\ru{вдова}}.) An heir? (\textbf{\ru{наследник}}.) A jubilee? (\textbf{\ru{юбилей}}.) An anniversary? (\textbf{\ru{годовщина}}.) A newborn? (\textbf{\ru{новорождённый}}.) A christening? (\textbf{\ru{крестины}}.) To christen? (\textbf{\ru{крестить}}.)
\end{tcolorbox}
